 \vspace{0.1cm}
\subsection{Hubble Tension}
An important aspect of knowing precise distance of celestial objects is the measurement of Hubble constant. In 1924, astronomer E. Hubble observed the extra-galactic Cepheids and estimated their distances using Leavitt's Law. He noticed that the farther objects are red-shifted which implies that they are receding away from us. He deduced a linear relation between the distances of galaxies and their recession velocity \cite{hubble1929relation}. This relation is known as \textit{Hubble law} which concludes the expansion of universe. The proportionality factor of the linear relation, $H_0$ is called the Hubble constant. The Hubble law expressed as:

\begin{equation}
	v = H_0 D
\end{equation}
 
 
 \begin{table}[h!]
 	\centering
 	\caption{Result of $H_0$ Key project of HST; \\ r(random error), s(systematic error) \\ Taken from \cite{freedman2001final} }
 	\begin{tabular}{c|c}
 		\textbf{Distance Method} & \textbf{Hubble constant (km/s/Mpc)} \\ \hline
 		Type I Supernovae & $71 \pm 2_r \pm 6_s$ \\
 		Tully-Fisher Realtion & $71 \pm 3_r \pm 7_s$ \\
 		Surface brighness fluctuation & $70 \pm 5_r \pm 6_s$ \\
 		Type II Supernovae & $72 \pm 9_r \pm 7_s$ \\
 		Fundametal Plane & $82 \pm 6_r \pm 9_s$ \\ \hline
 		\textbf{Combined Result} & $72 \pm 8 $
 	\end{tabular}
 \end{table} 
 
\noindent The precise estimation of the Hubble constant is one of the ambitious goals of modern astrophysics as it aids to determine the age of the Universe. The Hubble constant appears to lie between 67 and 73 (km/s/Mpc) as it approaches distinct values from different methods. The SH0ES project determined $H_0$ locally, using precise distance measure of Cepheids and SN Type Ia. According to SH0ES project, the $H_0$ approaches to 73.04 $\pm$ 1.04 km/s/Mpc \cite{riess2022comprehensive}. On the other hand, cosmological studies suggests a different value of $H_0$. For instance, Plank mission maps the cosmic microwave background anisotropies at angular resolution greater than 10 arcminute with $\Lambda CDM$ cosmological model. The estimated value of Hubble constant using Plack mission data is 67.4 $\pm$ 0.5 km/s/Mpc \citep{2020planck}. The deviation of 5$\sigma$ in the value of $H_0$ between two different methods challenging our understanding of nature. This discrepancy in estimating the exact value of $H_0$ is known as the Hubble Tension. 


\section{IRSB}
The basic idea behind this method is: the variation of angular diameter of targeted Cepheid is proportional to changing radius, $\Delta \theta = \frac{\Delta R}{d}$, where inverse of distance is a proportional constant. From the spectroscopic observation, Doppler shift in star's spectrum can be measured which gives the variation of radial velocity of the stellar surface. On integrating radial velocity over time, radius during the pulsation cycle can be estimated. Information about the angular diameter of targeted star can be deduced from interferometry technique by using the formula $\theta = \frac{lambda}{D}$, where $\lambda$ is observed wavelength and D is the baseline of interferometer. By substituting varying radius and angular diameter, distance of individual will be estimated. 


\subsection*{State-of-the-art Measurements}

The launch of the Gaia spacecraft in 2013 has revolutionized the field of astronomy. It measures the precise parallaxes and proper motion of celestial objects within the Milky Way and beyond, alongside, performs photometric observation in three bands (G, BP and RP) \cite{vallenari2022gaia} \cite{bailer2021estimating}. In June 2022, 34 months of observational data was published by Gaia collaboration - Gaia DR3, which provides astrometric data for $\sim$ 1.8 billion stars within the Milky Way. It is emerging as a revolution as the modern astronomers could push their research to a new leap by using Gaia parallax driven distances. 


In 2021, to confirm the Hubble tension with $\Lambda\mbox{CDM}$ model, Adam G. Riess used the Gaia parallaxes of 75 extragalactic Cepheids \cite{riess2021cosmic} to infer the zero-point of Leavitt Law. Then he calibrated the Gaia distances with Hubble space telescope-driven results and concluded that the Gaia parallaxes could provide distances up to 1\% precision. Alongside, he constructed a cosmic distance ladder (see fig. 1) with SNe Ia calibrators and estimated the Hubble constant as $H_0 = 73.2 \pm 1.3 \mbox{ km  s}^{-1} \mbox{Mpc}^{-1}$ with 4.2 $\sigma$ difference with Plank CMB observation under $\Lambda$CDM model.  


\begin{wrapfigure}{r}{0.5\textwidth}
	\begin{center}
		\vspace{-0.1cm}
		\includegraphics[width=7.5cm]{TRGB}
		\vspace{-0cm}
		\caption{A color magnitude diagram of 46 Galactic globular clusters depicting TRGB at $M_I = -4.056$ mag }
		\vspace{-0.5cm}
	\end{center}
\end{wrapfigure}
Another project to measure the Hubble constant from the TGBR method is done by W. Freedman. In this method, the brightest red giant stars - residing in the globular clusters; are used as a standard candle. It follows as the red giant stars burn Hydrogen in a shell and accumulate Helium ash in the core. Once enough Helium accumulates in the core as degenerate matter, at a certain pressure limit, it fuses and creates a spontaneous flash of immense energy, bringing the star down to the horizontal branch in color-magnitude diagram. Before the Helium flash, the red giant star touches the tip of the red giant branch upon which this method is named as TRGB method. Using this technique, \cite{freedman2021measurements} used the TRGB distances with Type Ia supernovae from Carnegie Supernova Project to form a distance ladder for deducing the value of the Hubble constant. She concluded the $H_0 = 69.8 \pm 0.6 \mbox{ (stat)} \pm 1.6 \mbox{ (sys)} \mbox{ km s}^{-1} \mbox{Mpc}^{-1}$.





\The golden data contains photometry with quality index less than or equal to 3 and RUWE less than 1.4 for Gaia parallax. Raw dataset contains 103 Galactic Cepheids pulsating in fundamental mode. Their periods are measured in days and scaled logarithmically. The distance modulus is derived from the Gaia DR3 parallax filtered with RUWE < 1.4, the color excess is adopted from Fernie measurements, and photometric magnitudes are collected in the BVIJHK bands. The distribution of each observable against the period is depicted in the following sequence of pair plots. A quick glance shows no clear correlation in the plots.




\subsubsection*{Target-Based Reddening Law}
The refined calibration methodology also provides insights into the reddening law for individual Cepheids. By comparing observed deviations across multiple bands and Wesenheit functions, star-specific reddening ratios can be inferred, allowing a more precise correction of extinction effects. This target-based approach can improve the accuracy of distance estimates, particularly in regions with non-standard or spatially varying extinction, and serves as a foundation for further refinement of both Galactic and extragalactic distance measurements.


\cite{2009MNRAS.396.1287S}


\subsection{Astrophysics, Astronomy, and Cosmology}

Astrophysics is inherently interdisciplinary, drawing upon principles from physics, chemistry, geology, computer science, and related fields to construct a coherent understanding of the Universe and humanity’s place within it. Astrophysicists investigate the physical processes governing celestial objects, their interactions and evolution, the dynamics of complex cosmic systems, and the potential existence of life beyond Earth.

Astronomy may be regarded as the historical precursor to astrophysics, with a primary emphasis on the observation and systematic description of celestial objects. It focuses on measuring positions, motions, brightness, and classifications, forming the empirical foundation upon which physical interpretations are built. This observational approach is essential for modeling the dynamics of the Universe and for predicting astronomical phenomena such as solar and lunar eclipses, planetary conjunctions, and stellar motions within the Galaxy, as well as for mapping the observable Universe. In this sense, astronomy primarily addresses questions of \emph{what} and \emph{where}, while astrophysics extends this framework to address questions of \emph{how} and \emph{why}.

On the largest scales, cosmology treats the Universe as a single, interconnected physical system. It seeks to understand the origin, geometry, age, evolution, and ultimate fate of the Universe as a whole. Across this hierarchy of spatial scales, galaxies emerge as the fundamental building blocks, which themselves are organized into larger structures such as groups, clusters, and superclusters. These structures are arranged along vast filamentary networks interspersed with low-density regions known as cosmic voids, together forming what is commonly referred to as the cosmic web.



Every new discovery in the history of science has served as a source of new insights. One such discovery was the identification of pulsating stars and the strong correlation between their pulsation period and luminosity, now known as the Leavitt Law \citep{1908leavitt}. This Master’s thesis is devoted to the subject of distance determination in astronomy, specifically focusing on the Leavitt Law — one of the primary methods used since the early twentieth century to determine extragalactic distances. This method relies on a special class of stars, Classical Cepheid variables, which pulsate radially and can be observed in distant galaxies due to their high luminosities. In this research, I utilize multiband photometric data of 121 Galactic Cepheids along with their measured distances and interstellar reddening along the line-of-sight. Using this comprehensive dataset, I demonstrate how multiband photometry can be employed to achieve high-precision estimates of distance and reddening, leading to a well-constrained Galactic Leavitt Law.

This introductory chapter provides a general background on the topic, tracing its historical development up to the present day. It also discusses how the discovery of the Leavitt Law fundamentally changed astronomers' perspectives on both the evolution and the scale of the Universe. Furthermore, this thesis is organized into three main parts, in addition to this introductory chapter. Part I comprises two chapters: the first details how key stellar parameters—such as luminosity, distance, and reddening—are derived from observational data; the second presents the physics of Cepheid pulsation and its relationship with the Leavitt Law. Part II also contains two chapters: the first explores the dataset available for calibration, and the second develops the calibration methodology to address systematic effects in distance and reddening. Part III consists of two chapters as well: the first compares the corrections suggested for two independent distance measurements (IRSB and Gaia parallax), and further verifies the calibrated distances of a sample of 17 Cepheids using an independent cluster-membership method while discussing the limitations of the current approach. The final chapter concludes with the calibrated Leavitt Law, presents its application to determine distances to the Magellanic Clouds and their three-dimensional geometries, and discusses prospects for future work, including data availability and potential extensions of the current algorithm.




\subsection{The Big Bang Model: 1930s}

Independently of Hubble’s observational work, Georges Lemaître applied Einstein’s equations of general relativity to cosmology and arrived at solutions describing an expanding Universe. In 1931, he proposed that cosmic expansion originated from an extremely hot and dense initial state, an idea that later evolved into the Big Bang model \citep{1931lemaitre}.

Further theoretical developments occurred in the late 1940s, when Gamow proposed a hot early Universe model, subsequently refined by Alpher and Herman \citep{1948gamow,1948alpher}. These works predicted the existence of a relic thermal radiation field—the cosmic microwave background (CMB)—as a remnant of the early Universe, with an expected temperature of a few kelvin.

This prediction was observationally confirmed in 1965, when Arno Penzias and Robert Wilson detected an isotropic microwave signal with a temperature of $3.4 \pm 1$~K using the Holmdel Horn antenna \citep{1965penzias}. Their discovery provided compelling empirical support for the Big Bang model, which today constitutes the foundation of the standard cosmological paradigm, describing the origin, evolution, and large-scale structure of the Universe.


\subsection{Technological Development: 1960s}

Despite persistent theoretical and observational challenges, technological progress has continually driven advances in astronomy. Beginning in the 1960s, the transition from photographic plates to electronic detectors marked a major transformation in observational capabilities. The introduction of charge-coupled devices (CCDs) in astronomical photometry dramatically improved sensitivity, observation range, and measurement precision \citep{2023roth}. While photographic plates could only detect 1–2 \% of incoming photons, CCDs can register up to 90\% of incident flux, enabling the observation of much fainter objects. An important advantage of CCD imaging is the digitalization of photometric data. Each CCD chip contains an array of photosensitive pixels, which emits electrons upon absorbing incoming photons via the photoelectric effect \citep{1905einstein}. Higher light intensity leads to higher voltage potential, corresponding to higher pixel values in the resulting area of image. The full grid of pixels digitally maps the segment of the sky exposed to the camera through telescope.

As the light of different frequencies interact with materials differently, CCD imaging technique enable astronomers to observer in wider range of wavelengths. Standard silicon-based CCDs operate efficiently within the optical band. Though, silicon CCDs coated with magnesium fluoride (\ce{MgF2}) extend the absorption in UV band \citep{2011hamden}. For Infrared band (700 - 2500 nm), mercury cadmium telluride (\ce{HgCdTe}) based detectors are used to achieve higher absorption\citep{2009fox}.

On the computational front, sustained advances in high-performance computing have profoundly transformed astronomical research by enabling detailed numerical investigations of complex astrophysical systems. These include simulations of stellar interiors \citep{1990kippenhahn}, galaxy formation and evolution \citep{1993navarro}, and cosmological structure formation within large-scale simulations \citep{2025TNGCluster}. Such studies rely on parallel computing and supercomputing architectures to solve coupled gravitational, hydrodynamical, and radiative processes across a wide range of spatial and temporal scales.

The rapid growth of data from ground- and space-based surveys has necessitated advanced infrastructures for data processing, archiving, and dissemination. Initiatives such as the Virtual Observatory (VO) provide standardized, interoperable frameworks for the efficient analysis of distributed astronomical datasets \citep{2009hatziminaoglou}. Complementing these efforts, specialized software tools such as TOPCAT enable efficient visualization, cross-matching, and statistical analysis of large astronomical catalogues while maintaining compatibility with VO standards \citep{2005TOPCAT}. Together, these tools underpin the efficiency, reproducibility, and scientific impact of modern data-driven astronomy.


Dark matter

Strong observational evidence for dark matter arises from gravitational lensing, galaxy rotation curves, and the fine-scale anisotropies observed in the cosmic microwave background (CMB). Importantly, the interpretation of these phenomena depends critically on accurate distance measurements; for example, lensing geometries and dynamical mass estimates are highly sensitive to the assumed distances of the relevant systems. 


\subsubsection{$\Lambda$CDM Cosmological Model}

Observations of the distant Universe allow us to probe earlier cosmic epochs, as light requires finite time to traverse space. This fundamental property enables direct investigation of the Universe’s evolution. One approach to characterizing cosmic expansion involves measuring distances to the most distant observable objects, effectively constraining the size of the observable Universe and the cosmological horizon.

By combining insights from high-energy particle physics with general relativity, cosmologists have been able to model the evolution of the Universe back to fractions of a second after the Big Bang. Integrating these theoretical frameworks with a broad range of observational data has led to the formulation of the standard cosmological model, known as the $\Lambda$CDM (Lambda Cold Dark Matter) model. This model is characterized by a small set of parameters, including the densities of dark energy, dark matter, and baryonic matter; the Hubble constant ($H_0$); the scalar spectral index describing primordial density fluctuations; and the curvature parameter that determines the geometry of spacetime. Many of these parameters are tightly constrained by measurements of the cosmic microwave background and are independently supported by large-scale structure surveys, baryon acoustic oscillations, and supernova observations.

The full-sky temperature anisotropy map shown on the introductory page is a snapshot of the early Universe obtained by the \emph{Planck} satellite, a mission of the European Space Agency operational between 2009 and 2013. This map depicts the cosmic microwave background radiation, the relic thermal emission from the epoch of recombination. The small temperature fluctuations observed across the sky correspond to primordial density perturbations in the early Universe. Over billions of years, these perturbations grew through gravitational instability, ultimately giving rise to galaxies, clusters, and the large-scale structure observed today.

The CMB represents the earliest observable relic of the Universe, providing a view of conditions approximately 380{,}000~years after the Big Bang. Its precise measurement enables robust constraints on fundamental cosmological parameters, including the age, composition, and expansion rate of the Universe. Among the most significant results derived from the \emph{Planck} mission is a high-precision determination of the Hubble constant, $H_0$, which plays a central role in understanding both the past evolution and the future fate of the Universe.


Applying this calibrated framework to the Hubble--Lemaître relation yields
\[
H_0 = 73.30 \pm 1.04 \; \mathrm{km\,s^{-1}\,Mpc^{-1}},
\]
which differs by $\sim 5\sigma$ from the value inferred by \emph{Planck} CMB observations under $\Lambda$CDM assumptions \citep{2022riess,2018planck}. This persistent discrepancy, known as the \emph{Hubble tension}, may reflect unaccounted systematic uncertainties or signal new physics beyond the standard cosmological model, and remains a central challenge in contemporary cosmology \citep{Verde2019, DiValentino2021}.


$\bar{W}^{12}_\lambda = \alpha^{12}_\lambda (\log P - 1) + \beta^{12}_\lambda$

$\bar{M}_\lambda = \alpha_\lambda (\log P - 1) + \beta_\lambda$

$\Delta W_\lambda^{12} = W_\lambda^{12} - \bar{W}^{12}_\lambda$

$\Delta M_\lambda = M_\lambda^{0} - \bar{M}_\lambda$


The vast scale of the Universe necessitates the use of specialized distance measures. Astronomical units are used within the Solar System, light-years or parsecs on stellar and Galactic scales, and parsec multiples for extragalactic distances, with redshift serving as a cosmological distance indicator in an expanding Universe. The relationships among these commonly used units are summarized in Table~\ref{table:distance}.

\begin{table}[h!]
	\centering
	\caption{Conversion of astronomical distance units}
	\begin{tabular}{c | c c c c} 
		 & Kilometer & Astronomical Unit & Light Year & Parsec \\  
		\hline
		Kilometer & 1 &  &  & \\ 
		Astronomical Unit & $1.49\times10^{8}$ & 1 &  & \\
		Light Year & $9.46\times10^{12}$ & $6.32\times10^{4}$ & 1 & \\
		Parsec & $3.08\times10^{13}$ & $2.06\times10^{5}$ & 3.26 & 1 \\
	\end{tabular}
	\label{table:distance}
\end{table}


